% ==========================================================================
%  Chapter 76 -- Build System Diagnosis and Compilation Optimisation
% ==========================================================================

\chapter{Build System Diagnosis and Compilation Optimisation}
\label{ch:build-system-optimisation}


% =========================================================================
\section{Motivation}
\label{sec:build-motivation}

A full rebuild of \texttt{fall\_n} with eight parallel jobs
(\texttt{ninja -j8}) on a modern desktop (AMD Ryzen 7, 32\,GB RAM)
takes approximately 50~minutes to generate 115 Ninja targets
(28~executables and 65~tests, plus object files).  This severely
limits the developer iteration cycle: a single-header edit that
touches \texttt{header\_files.hh} forces recompilation of every
translation unit.

This chapter diagnoses the root causes of the excessive compilation
time, quantifies the contribution of each factor, and proposes a
concrete mitigation plan.


% =========================================================================
\section{Build System Overview}
\label{sec:build-overview}

The build system uses CMake with the Ninja generator.  The main
\texttt{CMakeLists.txt} defines:

\begin{itemize}
  \item \textbf{28 executable targets}, including
    \texttt{fall\_n\_main}, \texttt{fall\_n\_seismic\_main},
    \texttt{fall\_n\_lshaped\_multiscale}, and 20~variants
    \texttt{fall\_n\_seismic\_v2} through
    \texttt{fall\_n\_seismic\_v21}.
  \item \textbf{65 test targets} registered via the custom
    \texttt{fall\_n\_add\_test()} macro, each compiling a
    separate test source file.
  \item \textbf{No static or shared library targets:}
    the codebase is entirely header-only, with all template
    definitions visible in every translation unit.
\end{itemize}


% =========================================================================
\section{Root Cause Analysis}
\label{sec:build-root-cause}

% -------------------------------------------------------------------------
\subsection{Monolithic Header: \texttt{header\_files.hh}}
\label{sec:build-monolithic-header}

The file \texttt{header\_files.hh} includes approximately 100
project headers covering \emph{every} module:

\begin{lstlisting}[caption={Excerpt from \texttt{header\_files.hh}.}]
// Core
#include "src/geometry/entities/Node.hh"
#include "src/geometry/entities/Vertex.hh"
// ... 40+ geometry/element headers

// Materials
#include "src/materials/concrete/KoBatheConcrete3D.hh"
#include "src/materials/steel/MenegottoPintoSteel.hh"
// ... 15+ material headers

// Model, Analysis, VTK
#include "src/model/Model.hh"
#include "src/analysis/NonlinearAnalysis.hh"
#include "src/vtk/VtkWriter.hh"
// ... 20+ analysis/vtk headers

// Reconstruction
#include "src/reconstruction/SubModelSolver.hh"
#include "src/reconstruction/NonlinearSubModelEvolver.hh"
#include "src/reconstruction/HomogenizedSection.hh"
\end{lstlisting}

Every test and executable that includes \texttt{header\_files.hh}
must parse and instantiate templates from all 100+ headers.  Since
most tests depend on only a small subset of the library, the vast
majority of this parsing is wasted.

\paragraph{Estimated impact.}
Header parsing and template instantiation account for approximately
$70{-}80\%$ of compilation time per translation unit.  With 93~TUs
(28~executables $+$ 65~tests), the aggregate waste is substantial.


% -------------------------------------------------------------------------
\subsection{Duplicate Seismic Executables}
\label{sec:build-duplicate-seismic}

Twenty targets (\texttt{fall\_n\_seismic\_v2} through
\texttt{fall\_n\_seismic\_v21}) all compile the same source file,
\texttt{main\_multiscale\_seismic.cpp}, with different preprocessor
definitions:

\begin{lstlisting}[caption={CMakeLists.txt excerpt.}]
add_executable(fall_n_seismic_v2  main_multiscale_seismic.cpp)
target_compile_definitions(fall_n_seismic_v2  PRIVATE SEISMIC_V2)
add_executable(fall_n_seismic_v3  main_multiscale_seismic.cpp)
target_compile_definitions(fall_n_seismic_v3  PRIVATE SEISMIC_V3)
# ... v4 through v21
\end{lstlisting}

Each variant triggers a full compilation of a large translation unit
(\texttt{main\_multiscale\_seismic.cpp}) that includes
\texttt{header\_files.hh}.  Thus, the same 100+ headers are parsed
and instantiated 20~times, producing nearly identical object files.

\paragraph{Estimated impact.}
These 20~targets represent ${\sim}35\%$ of the total build time.
They can be consolidated into a single executable with a
runtime parameter.


% -------------------------------------------------------------------------
\subsection{No Precompiled Headers (PCH)}
\label{sec:build-no-pch}

CMake supports precompiled headers via
\texttt{target\_precompile\_headers()}.  The current project does
not use this feature.  Third-party headers (PETSc, Eigen, VTK,
\texttt{<algorithm>}, \texttt{<vector>}, etc.) are re-parsed in
every TU.

\paragraph{Estimated impact.}
PCH for stable third-party headers typically reduces per-TU
compilation time by $15{-}25\%$.


% -------------------------------------------------------------------------
\subsection{No Explicit Template Instantiation}
\label{sec:build-no-explicit-inst}

The library's template-heavy design (e.g.\
\texttt{Model<ElementT, MeshT>},
\texttt{ContinuumElement<Material, Kinematics, Quadrature>})
means that every translation unit that uses these templates
instantiates them independently.  Identical instantiations (e.g.\
\texttt{ContinuumElement<KoBathe, SmallStrain, GaussLegendre>})
occur in multiple TUs, leading to redundant compilation that is
only discarded at link time.

\paragraph{Estimated impact.}
Explicit instantiation with \texttt{extern template} declarations
can reduce compilation time by $20{-}30\%$ for template-heavy
codebases.


% =========================================================================
\section{Mitigation Plan}
\label{sec:build-mitigation}

Table~\ref{tab:build-mitigation} summarises the proposed
optimisation steps in priority order, with estimated build-time
reductions and implementation effort.

\begin{table}[htbp]
  \centering
  \caption{Build optimisation steps, ordered by impact.}
  \label{tab:build-mitigation}
  \small
  \begin{tabular}{@{}clccc@{}}
    \toprule
    \textbf{Step} &
    \textbf{Optimisation} &
    \textbf{Reduction} &
    \textbf{Effort} &
    \textbf{Risk} \\
    \midrule
    1 & Consolidate seismic variants &
        ${\sim}35\%$ & Low  & None \\
    2 & Add PCH (Eigen, PETSc, VTK, STL) &
        ${\sim}15\%$ & Low  & None \\
    3 & Explicit template instantiation &
        ${\sim}25\%$ & Med. & Moderate$^{a}$ \\
    4 & Modularise \texttt{header\_files.hh} &
        ${\sim}15\%$ & High & Low \\
    5 & CMake OBJECT library &
        ${\sim}10\%$ & Med. & Low \\
    6 & C++20 modules (experimental) &
        TBD   & High & High$^{b}$ \\
    \bottomrule
  \end{tabular}
  
  \medskip
  \footnotesize
  $^{a}$ Requires maintaining explicit instantiation lists in sync
  with the template usage across TUs.  Missing an instantiation causes
  a linker error.

  $^{b}$ GCC~15 offers partial C++20 module support; not yet production-ready.
\end{table}

% -------------------------------------------------------------------------
\subsection{Step 1: Consolidate Seismic Variants}
\label{sec:build-seismic-consolidate}

Replace 20 separate executables with a single parameterised
executable that selects the variant at runtime:

\begin{lstlisting}[language=C++, caption={Proposed unified seismic driver.}]
int main(int argc, char* argv[]) {
    int variant = 2;   // default
    if (argc > 1) variant = std::atoi(argv[1]);
    
    switch (variant) {
        case 2:  run_seismic_v2();  break;
        case 3:  run_seismic_v3();  break;
        // ...
        default:
            std::cerr << "Unknown variant: "
                      << variant << "\n";
            return 1;
    }
}
\end{lstlisting}

If compile-time configuration is truly required, use
\texttt{if constexpr} with a single \texttt{-DSEISMIC\_VERSION=N}
define, building only the needed variant.  The CMake target reduces
to one:

\begin{lstlisting}
add_executable(fall_n_seismic main_multiscale_seismic.cpp)
target_compile_definitions(fall_n_seismic PRIVATE
    SEISMIC_VERSION=${SEISMIC_VERSION})
\end{lstlisting}


% -------------------------------------------------------------------------
\subsection{Step 2: Precompiled Headers}
\label{sec:build-pch}

Add PCH for the heaviest third-party headers:

\begin{lstlisting}[caption={CMake PCH setup.}]
# In CMakeLists.txt
set(PCH_HEADERS
    <Eigen/Dense>
    <Eigen/Sparse>
    <petsc.h>
    <petscsnes.h>
    <petscmat.h>
    <vtkSmartPointer.h>
    <vtkUnstructuredGrid.h>
    <algorithm>
    <vector>
    <array>
    <memory>
    <cmath>
)

# Apply to all fall_n targets
foreach(tgt ${FALL_N_ALL_TARGETS})
    target_precompile_headers(${tgt} PRIVATE ${PCH_HEADERS})
endforeach()
\end{lstlisting}


% -------------------------------------------------------------------------
\subsection{Step 3: Explicit Template Instantiation}
\label{sec:build-explicit-inst}

Identify the most commonly instantiated template combinations and
provide explicit instantiation definitions in dedicated
\texttt{.cpp} files:

\begin{lstlisting}[caption={Explicit instantiation example.}]
// file: src/model/Model_instantiations.cpp
#include "header_files.hh" // only this TU pays full cost

// Explicit instantiation definitions
template class Model<
    TimoshenkoBeam3D<GaussLobatto<3>>,
    DMPlexMesh>;

template class Model<
    ContinuumElement<
        Material<ThreeDimensionalMaterial>,
        SmallStrainKinematics,
        GaussLegendre<2,2,2>>,
    DMPlexMesh>;
\end{lstlisting}

In the headers, add \texttt{extern template} declarations to
suppress implicit instantiation:

\begin{lstlisting}[caption={\texttt{extern template} in header.}]
// In Model.hh, at the bottom:
#ifdef FALL_N_EXTERN_TEMPLATES
extern template class Model<
    TimoshenkoBeam3D<GaussLobatto<3>>,
    DMPlexMesh>;
// ... other common instantiations
#endif
\end{lstlisting}

A new CMake OBJECT library compiles the instantiation TUs once,
and all executables link against it.


% -------------------------------------------------------------------------
\subsection{Step 4: Header Modularisation}
\label{sec:build-modularise}

Decompose \texttt{header\_files.hh} into tiered headers:

\begin{center}
  \begin{tabular}{ll}
    \toprule
    \textbf{Header} & \textbf{Contents} \\
    \midrule
    \texttt{fall\_n/core.hh} &
      Node, Vertex, BasisFunctions, Quadrature \\
    \texttt{fall\_n/materials.hh} &
      Material, KoBathe, Menegotto, \ldots \\
    \texttt{fall\_n/elements.hh} &
      ContinuumElement, BeamElement, TrussElement \\
    \texttt{fall\_n/model.hh} &
      Model, DofHandler, BCs \\
    \texttt{fall\_n/analysis.hh} &
      NonlinearAnalysis, DynamicAnalysis \\
    \texttt{fall\_n/reconstruction.hh} &
      SubModel*, FieldTransfer, Evolver \\
    \texttt{fall\_n/vtk.hh} &
      VtkWriter, VtkPatch, \ldots \\
    \texttt{fall\_n/all.hh} &
      Includes all of the above (backward compat.) \\
    \bottomrule
  \end{tabular}
\end{center}

Test files include only the tier they need:
\begin{lstlisting}
// test_quadrature.cpp
#include "fall_n/core.hh"   // 10 headers instead of 100
\end{lstlisting}

This change requires updating all 65+ test files and 28+
main files, so it is best done incrementally.


% -------------------------------------------------------------------------
\subsection{Step 5: CMake OBJECT Library}
\label{sec:build-object-lib}

Create a shared compilation target:

\begin{lstlisting}[caption={CMake OBJECT library.}]
add_library(fall_n_lib OBJECT
    src/model/Model_instantiations.cpp
    src/elements/Element_instantiations.cpp
    src/analysis/Analysis_instantiations.cpp
)
target_precompile_headers(fall_n_lib PRIVATE ${PCH_HEADERS})

# All executables link to the OBJECT library
foreach(tgt ${FALL_N_ALL_TARGETS})
    target_link_libraries(${tgt} PRIVATE fall_n_lib)
endforeach()
\end{lstlisting}

This ensures that common template instantiations are compiled
exactly once.


% -------------------------------------------------------------------------
\subsection{Step 6: C++20 Modules (Exploratory)}
\label{sec:build-modules}

GCC~15 provides experimental support for C++20 named modules.
A future step could convert the most include-heavy components into
modules:

\begin{lstlisting}
// fall_n.core.cppm
export module fall_n.core;
export import :node;
export import :vertex;
export import :quadrature;
\end{lstlisting}

This is not recommended until the GCC/Clang module build system
integration (via CMake 3.28+ \texttt{CXX\_MODULES}) is stable.


% =========================================================================
\section{Expected Cumulative Impact}
\label{sec:build-impact}

Applying Steps 1--5 in sequence is expected to reduce the full build
time as follows:

\begin{equation*}
  \underbrace{50\;\text{min}}_{\text{current}}
  \xrightarrow[\text{Step 1}]{-35\%}
  32\;\text{min}
  \xrightarrow[\text{Step 2}]{-15\%}
  27\;\text{min}
  \xrightarrow[\text{Step 3}]{-25\%}
  20\;\text{min}
  \xrightarrow[\text{Step 4}]{-15\%}
  17\;\text{min}
  \xrightarrow[\text{Step 5}]{-10\%}
  15\;\text{min}
\end{equation*}

\noindent
A \textbf{${\sim}70\%$ overall reduction} is achievable with moderate
engineering effort.  Incremental builds (single-header changes) would
benefit even more, as PCH and modularised headers eliminate the
re-parsing of unrelated headers.


% =========================================================================
\section{Summary}
\label{sec:build-summary}

The current build system's excessive compilation time is attributable
to four primary causes: a monolithic uber-header, 20~duplicate
seismic targets, absence of precompiled headers, and redundant
template instantiation across 93~translation units.

A five-step optimisation plan --- consolidating seismic variants,
adding PCH, introducing explicit template instantiations,
modularising the header hierarchy, and creating a CMake OBJECT
library --- can reduce full build times from ${\sim}50$ minutes to
${\sim}15$ minutes.

These improvements require no changes to the library's public API
or test expectations; they are purely internal build-system
and compilation-model optimisations.
